\section{Proof of Work and Chain Selection}
\label{sec:pow}

Bitcoin blocks commit to proof of work through the block header hash. A block is
valid only if its header hash satisfies the target encoded by the header's
\code{nBits} field.

\subsection{Target encoding}

The \code{nBits} field is a 32-bit compact target. For compact value \(c\), let:

\[
  e = c \gg 24,\qquad
  m = c \mathbin{\&} 0x007fffff,\qquad
  s = c \mathbin{\&} 0x00800000.
\]

The decoded target \(\mathrm{C}^{-1}(c)\) is:

\[
\mathrm{C}^{-1}(c) =
\begin{cases}
  m \gg 8(3-e), & e \le 3,\\
  m \ll 8(e-3), & e > 3.
\end{cases}
\]

The sign bit is inherited from an older multiple-precision integer format and is
not permitted for proof-of-work targets. A decoded target is invalid if any of
the following conditions holds:

\begin{itemize}[leftmargin=*]
  \item \(s\) is set and \(m \ne 0\);
  \item \(\mathrm{C}^{-1}(c)=0\);
  \item decoding overflows 256 bits, which occurs when \(m \ne 0\) and either
        \(e > 34\), or \(m > 0xff\) and \(e > 33\), or \(m > 0xffff\) and \(e > 32\);
  \item \(\mathrm{C}^{-1}(c)\) is greater than the network's proof-of-work limit.
\end{itemize}

This compact-target decoding matches contemporary Bitcoin Core consensus
validation \cite{bitcoin-core-v31}.

The compact encoding \(\mathrm{C}(T)\) of target \(T\) is computed as follows. Let
\code{size = ceil(bitLength(T) / 8)}. If \code{size <= 3}, set
\code{word = T << 8(3 - size)}; otherwise set
\code{word = T >> 8(size - 3)}. If \code{word} has bit \code{0x00800000} set,
right-shift \code{word} by 8 and increment \code{size}. Then:

\[
  \mathrm{C}(T) = (size \ll 24) \mathbin{|} (word \mathbin{\&} 0x007fffff).
\]

Consensus proof-of-work targets are nonnegative, so the sign bit is not set in
valid block targets.

\subsection{Proof-of-work check}

For header \(h\), let:

\[
  T_h = \mathrm{C}^{-1}(h.\text{\code{bits}}),\qquad p_h = \mathrm{le}(\mathrm{bh}(h)).
\]

A candidate block satisfies proof of work iff \(T_h\) is a valid target and:

\[
  p_h \le T_h.
\]

A node rejects the block header if target decoding fails or if the hash integer
is greater than the target. The network-specific proof-of-work limit is part of
the chain parameters, not a value inferred from the genesis block
\cite{bitcoin-core-v31}.

\subsection{Difficulty adjustment}

Mainnet difficulty adjusts over fixed-height intervals. Test networks can have
different adjustment behavior. Network-specific exceptions must be cited from
chain parameter sources rather than inferred from mainnet behavior.

For mainnet:
\[
  \text{\code{DIFFICULTY\_ADJUSTMENT\_INTERVAL}} =
  \frac{\text{\code{TARGET\_TIMESPAN}}}{\text{\code{TARGET\_SPACING}}}.
\]
A candidate block's \code{nBits} must equal the value returned by the network's
difficulty-adjustment function for the previous block and candidate timestamp
\cite{bitcoin-core-v31}.

For networks without the minimum-difficulty exception, the algorithm is:

\begin{enumerate}[leftmargin=*]
  \item If the candidate height is not an exact difficulty-adjustment interval,
        return the previous block's \code{nBits}.
  \item Otherwise, identify the first block in the previous adjustment period.
        For mainnet this is the ancestor
        \code{DIFFICULTY\_ADJUSTMENT\_INTERVAL - 1} blocks before
        \code{previous}.
  \item Compute
        $\Delta = time(previous) - time(first)$.
  \item Clamp $\Delta$ to the range
        \[
          \frac{\text{\code{TARGET\_TIMESPAN}}}{\text{\code{MAX\_RETARGET\_FACTOR}}}
          \le \Delta \le
          \text{\code{TARGET\_TIMESPAN}} \cdot
          \text{\code{MAX\_RETARGET\_FACTOR}}.
        \]
  \item Decode the previous period's base target. On testnet4-like networks with
        BIP94 enforcement, the base target is taken from the first block of the
        period; otherwise it is taken from the previous block.
  \item Set:
        \[
          T' =
          \frac{baseTarget \cdot \Delta}{\text{\code{TARGET\_TIMESPAN}}}.
        \]
  \item If $T'$ exceeds the network proof-of-work limit, use the proof-of-work
        limit.
  \item Return \(\mathrm{C}(T')\).
\end{enumerate}

On no-retarget networks, the function returns the previous block's
\code{nBits}. On networks that allow special minimum-difficulty blocks, the
function may return the proof-of-work limit when the candidate timestamp is more
than \(2 \cdot \text{\code{TARGET\_SPACING}}\) after the previous block;
otherwise it searches backward for the most recent non-special difficulty in the
current adjustment period \cite{bitcoin-core-v31}.

\subsection{Accumulated work}

Nodes compare competing valid chains by accumulated work rather than by block
count. This follows the proof-of-work security model described in the Bitcoin
whitepaper \cite{nakamoto-bitcoin}.

For each accepted header, the node adds the work implied by that header's target
to the predecessor's accumulated work. Headers and blocks with invalid proof of
work are rejected before they can contribute to active-chain selection.

For a valid compact target, Bitcoin Core computes the work represented by a
header as:

\[
  work(c) =
  \left\lfloor \frac{2^{256}-\mathrm{C}^{-1}(c)-1}{\mathrm{C}^{-1}(c)+1} \right\rfloor + 1.
\]

This is algebraically equivalent to \(\lfloor 2^{256}/(\mathrm{C}^{-1}(c)+1) \rfloor\) while
avoiding construction of the unrepresentable integer $2^{256}$ in a 256-bit
integer type \cite{bitcoin-core-v31}. The active chain is selected from
valid candidates by greatest accumulated work, not by height.
